\documentclass[11pt]{article}
\usepackage[letterpaper,margin=1in]{geometry}
\usepackage[T1]{fontenc}
\usepackage{lmodern,microtype,graphicx,booktabs,amsmath,amssymb,array,hyperref,fancyhdr,caption}
\usepackage[section]{placeins}
\hypersetup{hidelinks,pdfdisplaydoctitle=true,pdftitle={HPS Gaussian-Process Resonance Search: Analysis Note v5.8.5},pdfauthor={Emrys Peets},pdfsubject={Common-mass likelihood and mass-coherence extension}}
\urlstyle{same}
\pagestyle{fancy}\fancyhf{}
\setlength{\headheight}{14pt}
\fancyhead[L]{\small HPS Gaussian-Process Resonance Search}
\fancyhead[R]{\small v5.8.5 / Review Draft}
\fancyfoot[C]{\thepage}
\fancypagestyle{plain}{\fancyhf{}\fancyhead[L]{\small HPS Gaussian-Process Resonance Search}\fancyhead[R]{\small v5.8.5 / Review Draft}\fancyfoot[C]{\thepage}}
\setlength{\tabcolsep}{5pt}
\setlength{\emergencystretch}{2em}
\renewcommand{\arraystretch}{1.12}
\captionsetup{font=small,labelfont=normalfont}
\renewcommand{\topfraction}{.95}\renewcommand{\bottomfraction}{.85}\renewcommand{\textfraction}{.05}\renewcommand{\floatpagefraction}{.8}\setcounter{topnumber}{3}\setcounter{bottomnumber}{2}\setcounter{totalnumber}{5}
\newcommand{\heading}[1]{\section{#1}}
\newcommand{\fig}[3][1]{\begin{figure}[htbp]\centering\includegraphics[width=#1\linewidth,height=.68\textheight,keepaspectratio]{../figures/#2.pdf}\caption{#3}\label{fig:#2}\end{figure}}
\newcommand{\sfN}{\overline\Phi}
\title{HPS Gaussian-Process Resonance Search\\[0.5em]
\large Common-mass likelihood and mass-coherence study\\Analysis Note v5.8.5}
\author{Emrys Peets\\Stanford University and SLAC National Accelerator Laboratory}
\date{19 September 2026}
\begin{document}
\maketitle
\begin{abstract}
We examine the structure near 92 MeV in the HPS 2015, 2016, and 2021
(10\%) samples using a common-mass resonance likelihood with independent
nonnegative signal amplitudes. The scan covers 19--250 MeV in 1 MeV steps,
with each dataset included within its declared support. For the nominal
blind half-width of $2.25\sigma_m$, the combined maximum occurs at 92 MeV.
The conditional Gaussian-response calibration gives local and full-domain
significances of $3.583$ and $2.060$, respectively. Seven of 256 direct
Poisson scans exceed the observed maximum, giving a 95\% probability
interval of $[0.0111,0.0555]$. Across four blind widths, the combined
maximum remains at 91--92 MeV. Individual peak stability is less uniform:
the 2021 diagnostic maximum switches from 93 MeV to a competing excursion
at 80 MeV. A toy-calibrated mass-coherence statistic gives $p=0.763$,
showing no unusual clustering under the specified null. These probabilities
condition on the pinned GP background source and exclude corrections for
choices made after inspecting the data.

\end{abstract}
\section{Introduction}
This continuation tests a resonance at one common mass with independent nonnegative signal amplitudes in the 2015, 2016 and 2021 (10\%) samples. It reuses the pinned v5.8.4 spectra, likelihoods, background sources, response derivatives and 256 complete Poisson scans. The full search remains the integer grid 19--250 MeV, with each dataset included only on its inherited support.
\begin{table}[htbp]\centering{\small\setlength{\tabcolsep}{4pt}\begin{tabular}{lrrrrr}
\toprule
Method & Peak [MeV] & Local $Z$ & Global $p$ & Global $Z$ & Direct $k/N$ \\
\midrule
Shared coupling & 66 & 2.83 & 0.2085 & 0.81 & 41/256 \\
Fisher & 92 & 3.72 & 0.0113 & 2.28 & 4/256 \\
signed Stouffer & 92 & 3.93 & 0.0049 & 2.58 & 1/256 \\
Free amplitudes & 92 & 3.58 & 0.0197 & 2.06 & 7/256 \\
\bottomrule
\end{tabular}}
\caption{Full-domain peaks at $\pm2.25\sigma_m$. Local and global $Z$ use the conditional Gaussian-response model. Direct counts use the same 256 full Poisson realizations. Old methods are unchanged v5.8.4 results; the new test has its own calibrated null distribution.}
\label{tab:baseline_table}
\end{table}

\fig[1]{methods_full_domain}{Baseline $\pm2.25\sigma_m$ comparison: local significance and probability for an equally strong score anywhere on the full 19--250 MeV grid. Curves use conditional Gaussian-response calibration; Table 1 retains the direct Poisson count for the new test. Shading marks the post-hoc 80--105 MeV diagnostic interval. Each method uses its own declared ordering and null maxima.}


\heading{Local and domain-global probabilities}
At a mass fixed independently of the data, the local probability is
\begin{equation}
 p_{\rm loc}(m)=\Pr_{0}\{T(m)\geq T_{\rm obs}(m)\},
 \qquad Z_{\rm loc}(m)=\Phi^{-1}[1-p_{\rm loc}(m)],
\end{equation}
where larger $T$ is more signal-like and the probability includes equality.
This inclusive convention matters when a nonnegative fitted amplitude gives
a point mass at $T=0$: its inclusive tail is one. A displayed nonnegative
$Z$ is a plotting convention; the probability remains the primary quantity.
The common-mass alternative has one amplitude for each active dataset, so
$\sqrt{q_R}$ is not generally its local significance. Its null distribution
must be calibrated with the appropriate boundary and active dataset count.

For a declared mass grid $\mathcal G$ and a declared local ordering $S(m)$,
the scan statistic and its domain-global probability are
\begin{equation}
 M=\max_{m\in\mathcal G} S(m),\qquad
 p_{\mathcal G}=\Pr_0\{M\geq M_{\rm obs}\}.
\end{equation}
For example, $S=-\log p_{\rm loc}$ compares locally calibrated evidence
across masses; $S=q_R$ defines a different valid scan test. The observed
maximum and every null maximum must use the same ordering. A full-spectrum
toy supplies the entire correlated curve, rather than independent random
fluctuations at each mass. This accounts for the look-elsewhere effect within
that grid. See
\href{https://arxiv.org/abs/1005.1891}{Gross and Vitells} for scan inference and
\href{https://arxiv.org/abs/2206.12328}{Ananiev and Read} for correlated Gaussian
significance-field approximations.

Here the inherited full combined grid is $19,20,\ldots,250$ MeV. The active
supports are 19--100 MeV (2015), 39--180 MeV (2016), and 50--250 MeV
(2021, 10\%); only 50--100 MeV contains all three datasets. ``Common mass''
ties the mass of the active components and does not extend an experiment's
support. Adding sites to a domain cannot reduce its null exceedance
probability at a \emph{fixed threshold}. Comparisons of observed maxima can
also change the threshold. The number of grid points is not the number of
independent trials: nearby masses share events, background fits, and signal
templates. Finer sampling resolves maxima until grid convergence; a
1-MeV-grid calibration is not, by itself, a continuous-mass calibration.

The 80--105 MeV interval was chosen after inspection and is used only to
describe the reported structure. Its 2015 interval ends at 100 MeV.
Calibrating a statistic restricted to this interval answers a conditional
diagnostic question; it does not replace the full-domain discovery test or
correct the choice of this interval. Likewise, four width-specific global
probabilities are four conditional results, not a selection-corrected claim
based on the most favorable width. Such a claim would require repeating the
entire width/method/domain selection in joint null realizations. Reusing
the same fluctuations across widths preserves their strong dependence.

All probabilities here condition on the pinned observed-data-derived GP
source, its kernel policy, and the stated likelihood construction. Direct
Poisson scans repeat the local fits; Gaussian response fields approximate
their fluctuations. Neither propagates uncertainty in choosing or estimating
the source itself. Exact likelihood factorization therefore does not imply
unconditional physical-discovery calibration. Boundary asymptotics are
reviewed by \href{https://arxiv.org/abs/1007.1727}{Cowan et al.}; sparse toy
tails here are reported with exceedance counts and binomial bounds.



\heading{Common-mass likelihood}
For independent datasets, at fixed mass and width,
\[
 L(\boldsymbol\mu,\boldsymbol\theta;m,w)
 =\prod_{d\in A(m)}L_d(\mu_d,\theta_d;m,w),\qquad \mu_d\geq0.
\]
Each local likelihood is the archived Poisson model with correlated Gaussian GP-background constraint. Its expected bin counts have the form
$\lambda_d=b_d+L_d\theta_d+A_d S_d$ and penalty $\theta_d^T\theta_d/2$. The engine uses the original template normalization $A_d$, with fitted-window yield $\mu_d=A_d\sum_i S_{di}$. This is a one-to-one amplitude reparameterization. Here $L_d$ in the count equation is the covariance factor, distinct from the likelihood symbol. The local background prediction and covariance are conditioned on the sidebands of each observed or toy spectrum. The full generating source is fixed.
\[
 q_R(m;w)=2\left[\sum_d\mathrm{NLL}_d(0,\hat\theta_{d,0})
       -\sum_d\mathrm{NLL}_d(\hat\mu_d^{+},\hat\theta_d)\right]
       =\sum_{d\in A(m)}[\max(r_d(m;w),0)]^2.
\]
The archived signed root $r_d$ uses the unconstrained single-amplitude fit and the zero-amplitude profile. Convexity and factorization make the clipped-root identity exact for this likelihood, including boundary fits. This is a direct likelihood construction; Fisher or Stouffer p-values do not enter $q_R$. Independent amplitudes impose no common coupling or calibrated relative yield. They also allow one or two amplitudes to vanish, so rejection of the all-zero null does not establish a signal in every dataset.

The exact likelihood statistic must be distinguished from its sampling calibration. For the conditional response approximation,
\[
 r_d^*(m,w)=a_d(m,w)+D_{d,m,w}^{T}\xi_d,\qquad
 \xi_d\sim N(0,I),\qquad s_d^2=D_{d,m,w}^{T}D_{d,m,w}.
\]
Different datasets fluctuate independently. Within each dataset the same latent fluctuation drives the entire mass scan. The probability of $q_R=0$ is
$\prod_{d\in A(m)}\Phi[-a_d(m,w)/s_d(m,w)]$; its inclusive tail is one.
The local tail of the sum of clipped squared normals is integrated over all nonempty positive-amplitude subsets with deterministic polar quadrature. At $a_d=0,s_d=1$ it reduces to the independent-boundary mixture $2^{-k}\sum_{j=1}^{k}\binom{k}{j}\Pr(\chi_j^2\geq q)$ for $q>0$, not a single $\chi_k^2$ tail. Actual mass-dependent $a_d,s_d$ are retained. Monotone lookup tables accelerate the scans. The declared ordering is $S_R(m)=-\log p_R(m)$.

The 256 stored Poisson realizations give an additional direct local and global calibration of this ordering, using the exact profile roots at every mass. At 92 MeV none exceeds the observed local statistic for any width; the one-sided 95\% upper bound is $p<0.01163$, which cannot resolve the response-model estimate near $3.6\sigma$. Zero counts are never presented as zero probability. Point estimates elsewhere use $(k+1)/(N+1)$; two-sided 95\% intervals are Clopper--Pearson. A separate 100,000-draw response ensemble preserves mass and width correlations through the joint derivative covariance.

\fig[1]{profile_statistic}{The exact observed profile statistic $q_R$ for all four widths. Left: full domain, including changes in active datasets. Right: diagnostic interval. The global ordering uses calibrated $-\log p_R$, whose maximum can differ from the maximum of raw $q_R$. Neither ordinate is replaced by $\sqrt{q_R}$ as a significance.}


\heading{Results of the combined search}
\begin{table}[htbp]\centering{\small\setlength{\tabcolsep}{4pt}\begin{tabular}{rrrrrrl}
\toprule
$w$ & Peak [MeV] & Local $Z$ & $Z(92)$ & Global $Z$ & $k/256$ & Direct global $p$: 95\% CI \\
\midrule
2.25 & 92 & 3.583 & 3.583 & 2.060 & 7 & [0.0111, 0.0555] \\
2.4 & 91 & 3.518 & 3.510 & 1.960 & 8 & [0.0136, 0.0606] \\
2.5 & 92 & 3.551 & 3.551 & 2.017 & 6 & [0.0086, 0.0503] \\
2.6 & 92 & 3.610 & 3.610 & 2.104 & 6 & [0.0086, 0.0503] \\
\bottomrule
\end{tabular}}
\caption{Free-amplitude peaks selected by maximum $-\log p_R$ on the full grid. The direct count column gives $k$ out of 256. Significance columns are Gaussian-response estimates. Raw $q_R$ instead has its maximum at 91 MeV for all four widths.}
\label{tab:width_table}
\end{table}

\fig[.94]{free_widths_diagnostic}{Free-amplitude local significance and full-domain probability, shown near the inspected structure. Every global probability in this zoom still uses the complete 19--250 MeV scan. Gaussian-response estimates are accompanied by direct Poisson counts and intervals in Table 2 and the numerical ledger. The four curves reuse the same observations and are not four independent confirmations.}
At 92 MeV, $q_R=17.307$--18.077, local $Z=3.510$--3.610, and full-domain Gaussian-response $Z=1.949$--2.104. At each width's selected peak, the direct global plus-one probabilities are 0.0311, 0.0350, 0.0272, 0.0272, corresponding to $Z$ values 1.864, 1.812, 1.923, 1.923. Their intervals are much wider than the 100,000-field Monte Carlo errors. The new test is less significant than signed Stouffer: it tests a broader alternative with separately nonnegative amplitudes, and its own boundary-aware local calibration.

The same width controls both the bins excluded from GP conditioning and the fitted likelihood window. Signal templates retain their original normalization. These are fixed-kernel window comparisons; no width has been adopted by choosing its largest observed significance. The displayed mass-global probabilities do not correct the post-observation comparison of widths or methods.


\heading{Mass stability across blind widths}
\fig[.98]{individual_stability}{All twelve signed reference-score curves $Z_{d,w}(m)=(r_{d,w}-a_{d,w})/s_{d,w}$ in the post-hoc diagnostic interval. Dots mark the 1 MeV grid; the dotted vertical line is 92 MeV. Gray space beyond 100 MeV in the 2015 panel is outside its support. Signed scores retain deficits, unlike nonnegative significance displays.}
\begin{table}[htbp]\centering{\small\setlength{\tabcolsep}{4pt}\begin{tabular}{lrrrrrrr}
\toprule
Dataset & $2.25$ & $2.4$ & $2.5$ & $2.6$ & $\sigma(92)$ & $\Delta m$ & $R$ \\
\midrule
2015 & 92 & 93 & 93 & 93 & 4.804 & 1 & 0.208 \\
2016 & 91 & 91 & 91 & 91 & 3.788 & 0 & 0.000 \\
2021 & 93 & 93 & 80 & 80 & 2.449 & 13 & 5.309 \\
\bottomrule
\end{tabular}}
\caption{Diagnostic-region signed-score argmax masses [MeV] at each half-width; resolution and range are also in MeV. $\Delta m=\max_w\hat m-\min_w\hat m$ and $R=\Delta m/\sigma(92)$. The region is 80--105 MeV, clipped to 80--100 for 2015. All selected observed peaks also pass the positive-raw-fit gate used by the coherence test.}
\label{tab:stability_table}
\end{table}

The 2015 and 2016 maxima remain within one grid bin across width. In 2021 the first two widths have almost tied excursions at 80 and 93 MeV: at $w=2.25$, their signed scores are 1.6216 and 1.6228. At $w=2.5$ they are 1.7332 and 1.6194. Thus the 13 MeV range is an argmax rank switch involving the lower diagnostic boundary, not evidence that a tracked resonance shifts by 13 MeV. The response at 92 MeV remains positive but modest.



\subsection{Individual amplitudes at 92 MeV}
\fig{stability_matrix}{Signed-score stability matrix: rows are dataset and half-width, columns are tested masses. The gray cells are unavailable 2015 masses. This is a post-hoc diagnostic display, with one common color scale; it is not a reduced discovery search.}
\begin{table}[htbp]\centering{\small\setlength{\tabcolsep}{4pt}\begin{tabular}{lrrrrr}
\toprule
Dataset & $w$ & Signed $Z(92)$ & Local $p(92)$ & $\hat\mu_d^+$ [events] & $\hat A_d$ \\
\midrule
2015 & 2.25 & 2.594 & 0.0047 & 1\,147 & 7748.35 \\
2015 & 2.4 & 2.594 & 0.0047 & 1\,173 & 7857.10 \\
2015 & 2.5 & 2.596 & 0.0047 & 1\,188 & 7922.75 \\
2015 & 2.6 & 2.691 & 0.0036 & 1\,233 & 8196.43 \\
2016 & 2.25 & 2.796 & 0.0026 & 16\,252 & 913.25 \\
2016 & 2.4 & 2.726 & 0.0032 & 16\,182 & 902.53 \\
2016 & 2.5 & 2.778 & 0.0027 & 16\,662 & 924.29 \\
2016 & 2.6 & 2.770 & 0.0028 & 16\,703 & 924.58 \\
2021 & 2.25 & 1.423 & 0.0774 & 7\,585 & 165.47 \\
2021 & 2.4 & 1.363 & 0.0864 & 7\,436 & 159.77 \\
2021 & 2.5 & 1.363 & 0.0864 & 7\,436 & 159.77 \\
2021 & 2.6 & 1.378 & 0.0841 & 7\,650 & 163.06 \\
\bottomrule
\end{tabular}}
\caption{Conditional single-dataset local responses and fitted amplitudes at the inspected mass 92 MeV. Local tails use the source-standardized response model. The event yield is integrated over the actual fitted window; changing the window changes its contained fraction.}
\label{tab:amplitude_table}
\end{table}

Here $\hat\mu_d^+$ is the independent nonnegative signal amplitude expressed as fitted signal events in that width's likelihood window. These are conditional fitted yields; they are not efficiency-corrected rates. The full-template yield and the original amplitude parameter $A_d$ are also saved in the ledger. In that parameterization $A_d=\epsilon_d^2/10^{-8}$, used here simply as an independent template normalization for each dataset. All twelve 92 MeV fitted amplitudes are positive. Their local reference $Z$ values alone do not demonstrate equal peak locations or a common physical coupling.


\heading{Null distribution of the coherence statistic}
For comparison on identical support, the coherence statistic uses the post-hoc common interval 80--100 MeV. In each dataset and width, select the maximum signed reference score among sites with a positive raw fit. Exact ties choose the lower mass. If any dataset--width has no positive site, its complete diagnostic is assigned $T=W=+\infty$ and cannot count as coherent.
\[
 \bar m=\frac{\sum_{d,w}\hat m_{d,w}/\sigma_d^2(92)}{\sum_{d,w}1/\sigma_d^2(92)},\qquad
 T=\frac1{12}\sum_{d,w}\left[\frac{\hat m_{d,w}-\bar m}{\sigma_d(92)}\right]^2,
 \qquad W=\frac13\sum_d R_d^2.
\]
Small $T$ describes peaks clustered at one mass across both datasets and widths. $W$ separately describes stability across width. The detector resolutions supply a dimensionless comparison scale; they are not claimed to be peak-estimator standard errors. These definitions were fixed in the continuation before extracting its observed peak table, but the region and scientific question arose after inspection. Their lower-tail probabilities are exploratory null diagnostics, not discovery significances, and are not multiplied together.
\begin{table}[htbp]\centering{\small\setlength{\tabcolsep}{4pt}\begin{tabular}{lrrrlr}
\toprule
Statistic & Observed & Direct $k/N$ & Direct $p$ & Direct 95\% interval & Gaussian $p$ \\
\midrule
T & 2.9793 & 195/256 & 0.7626 & [0.7047, 0.8126] & 0.7387 \\
W & 9.4102 & 253/256 & 0.9883 & [0.9661, 0.9976] & 0.9858 \\
\bottomrule
\end{tabular}}
\caption{Inclusive lower-tail probabilities. $T$ is the primary coherence diagnostic; $W$ is a separate width-stability diagnostic. The region and statistics have not received a selection correction.}
\label{tab:coherence_table}
\end{table}

\fig[.95]{coherence_null}{Inclusive lower-tail distributions from the 256 paired exact Poisson scans and 100,000 supplemental Gaussian-response scans. Each toy repeats the peak selection on the same grid and the same widths. Dashed lines mark observed values; smaller values indicate tighter coherence or greater stability. Ineligible toys remain in the denominator.}
Under the pinned null, 195 of 256 exact scans are at least as tightly clustered as the observed twelve peaks. This is not unusually coherent. Likewise, 253 of 256 have width ranges at least as small. The Gaussian estimates agree within the direct intervals; increasing Gaussian statistics cannot remove source-model uncertainty. The paired Gaussian ensemble contains 63 ineligible scans, all retained in the denominator as noncoherent. No diagnostic tail is used to select a new mass region or nominal width.



\heading{Validation}
The numerical package contains 928 free-amplitude mass--width rows and 12 independent 92 MeV amplitude checks. The joint constrained likelihood was replayed at 20 representative mass--width configurations, including single-, two- and three-dataset regions, and checked against the factorized statistic and nested shared-coupling fit. All source Poisson spectra regenerate exactly from their recorded seeds, and all twelve observed 92 MeV signed roots reproduce the parent values exactly in the matching SciPy/NumPy environment.

The central independent-boundary mixture and the positive/zero-atom limits pass. Increasing quadrature order changes the checked tails by at most 1.6e-14 relatively; lookup interpolation changes observed $\log p$ by at most 2.4e-07. Joint covariance reconstruction differs by at most 7.6e-12. Direct scalar evaluation reproduces the coherence metric to 4.4e-16. Numerical checks establish implementation consistency, not physical source validity.

At each width's Gaussian 95th-percentile global threshold, the direct Poisson exceedance counts are 16, 17, 20 and 17 out of 256. These are above the nominal mean 12.8; all corresponding 95\% binomial intervals still contain 0.05. This limited ensemble cannot rule out a modest response-tail mismatch. At the observed peaks, all Gaussian global estimates lie inside the direct intervals.

\heading{Discussion}
The free-amplitude likelihood identifies a combined excess at 91--92 MeV
throughout the four width choices. At the nominal width, the individual
local reference significances at 92 MeV are $2.59$, $2.80$, and $1.42$ for
2015, 2016, and 2021, respectively. The combined local significance is
$3.583$; scanning the full declared domain reduces it to $2.060$ in the
Gaussian-response approximation. Direct Poisson scans give a broader
assessment of this tail, with seven exceedances among 256 scans.

The stable combined maximum does not extend to every individual regional
maximum. In particular, a competing 2021 excursion overtakes the 93 MeV
peak at the larger widths. The mass-coherence probability of $0.763$
therefore provides no additional evidence for unusually aligned peaks.
Both conclusions remain conditional on the fitted background source and
the diagnostic choices prompted by the observed spectra.

The background source was fitted to the observed data and then frozen. Neither ensemble propagates source-estimation uncertainty, possible signal absorption during source construction, background-model selection, or unspecified correlated systematic uncertainties. The direct toys do refit the local GP predictions and profile the likelihood on each complete generated spectrum. The response fields approximate those operations by their stored derivatives. Agreement with 256 toys is a limited check, especially in sparse tails, not proof of precise Gaussian-tail calibration.

The 2015 support extension through 100 MeV, 2021 sample fraction, resolution functions, kernel states and signal templates remain inherited assumptions. No continuous-mass convergence claim is made: these results describe the 1 MeV grid. The parent study separately examined 0.5 versus 1 MeV spacing. The current extension does not choose a finer grid or smaller search domain after seeing a preferable result.


\subsection{Reproducibility}
The source bundle includes all pinned inputs needed by this continuation, copied likelihood code, analysis and plotting scripts, the protocol, numerical ledgers, saved null maxima and coherence statistics, validation records, and SHA-256 manifests. The released v5.8.4 study and output package were hashed before and after the continuation. Rebuild commands are in \texttt{README.md}; the original report is included unchanged and precedes this extension in the combined PDF.
\FloatBarrier
\begin{thebibliography}{9}
\bibitem{cowan} G. Cowan, K. Cranmer, E. Gross and O. Vitells,
\textit{Asymptotic formulae for likelihood-based tests of new physics},
\href{https://arxiv.org/abs/1007.1727}{arXiv:1007.1727}.
\bibitem{gross} E. Gross and O. Vitells,
\textit{Trial factors for the look elsewhere effect in high energy physics},
\href{https://arxiv.org/abs/1005.1891}{arXiv:1005.1891}.
\bibitem{ananiev} V. Ananiev and A. L. Read,
\textit{Gaussian Process-based calculation of look-elsewhere trials factor},
\href{https://arxiv.org/abs/2206.12328}{arXiv:2206.12328}.
\bibitem{parent} HPS-GPR v5.8.4,
\textit{Independent combinations, mass spacing and blind-window comparisons},
pinned report and inputs, 18 September 2026.
\end{thebibliography}
\end{document}
